\documentclass[12pt]{article}

% ---------------------------------------------------------
% PACKAGES
% ---------------------------------------------------------
\usepackage[a4paper,margin=0.8in,headheight=14.49998pt]{geometry}
\usepackage{hyperref}
\usepackage{titlesec}
\usepackage{setspace}
\usepackage{enumitem}
\usepackage{microtype}
\usepackage{fancyhdr}
\usepackage{tabularx}
\usepackage{seqsplit} % For wrapping long hashes
\usepackage{newtxtext,newtxmath}

% ---------------------------------------------------------
% DOCUMENT SETTINGS
% ---------------------------------------------------------
\hypersetup{
    colorlinks=true,
    linkcolor=black,
    urlcolor=blue,
    citecolor=black,
    pdfauthor={Leonard Okyere Afeke},
    pdftitle={Phishing Incident Report – Mutawa Marine Services Impersonation}
}

\setstretch{1.2}
\sloppy
\setlength{\emergencystretch}{3em}

\titleformat{\section}{\large\bfseries}{\thesection}{1em}{}
\titleformat{\subsection}{\normalsize\bfseries}{\thesubsection}{1em}{}

\pagestyle{fancy}
\fancyhf{}
\rhead{Phishing CTI Report}
\lhead{Leonard Okyere Afeke}
\cfoot{\thepage}

% ---------------------------------------------------------
% TITLE PAGE
% ---------------------------------------------------------
\begin{document}

\begin{titlepage}
    \centering
    \vspace*{3cm}
    {\Huge \textbf{Phishing Incident Report}}\par
    \vspace{0.5cm}
    {\Large Mutawa Marine Services Impersonation}\par
    \vspace{1.5cm}
    
    {\large \textbf{Author:} Leonard Okyere Afeke}\par
    \vspace{0.3cm}
    {\large \textbf{Date:} \today}\par
    \vspace{2cm}
    
    \vfill
    {\small This report is based entirely on publicly available information (OSINT). 
    No scanning, exploitation, or interaction with live systems occurred.}
\end{titlepage}

% ---------------------------------------------------------
% EXECUTIVE SUMMARY
% ---------------------------------------------------------
\section*{Executive Summary}

A phishing email impersonating Mutawa Marine Services was delivered to the recipient with a malicious RAR archive disguised as a PDF receipt. The attachment contained a Trojan Downloader, confirmed by VirusTotal (46/63 detections) and Hybrid Analysis (Threat Score 100/100). The email did not originate from the legitimate organization and instead came from a Hostwinds VPS, indicating clear domain spoofing.

The malware exhibited outbound encrypted communication (SSL/TLS) and TOR-related behavior, consistent with downloader malware attempting to contact remote infrastructure for follow-on payload delivery. VirusTotal’s Relations graph linked the payload to multiple phishing emails and ZIP archives across 2025, suggesting the sample is part of a broader commodity malware distribution campaign rather than a targeted attack.

No evidence indicates that the threat actor had specific interest in the recipient. The infrastructure, malware type, and delivery method all align with low-sophistication cybercrime activity aimed at establishing initial access for secondary payload deployment such as credential stealers, RATs, or ransomware.

% ---------------------------------------------------------
% INCIDENT NARRATIVE
% ---------------------------------------------------------
\section{Incident Narrative}

\subsection{Initial Delivery}
The phishing email impersonated Mutawa Marine Services and claimed to provide a shipping receipt. The attachment, named to resemble a PDF, was actually a RAR archive containing a disguised executable. This social engineering approach is consistent with commodity phishing campaigns that rely on user interaction to initiate malware execution.

\subsection{Email Infrastructure Analysis}
Header analysis revealed that the email did not originate from the legitimate mutawamarine.com domain. Instead, it was sent from a Hostwinds VPS (192.119.71.157), a low-cost hosting provider frequently used for disposable phishing infrastructure.

Authentication checks showed:
\begin{itemize}[noitemsep]
    \item SPF: Fail
    \item DKIM: None
    \item DMARC: Fail
\end{itemize}

The Reply-To address pointed to a mail.com account, a tactic commonly used to redirect victim responses away from the impersonated organization.

\subsection{Malware Analysis}
The attached archive (\texttt{SWT\_ \#09674321\_\_\_\_PDF\_\_.CAB}) contained a .NET executable (\texttt{SWT\_ \#09674321\_\_PDF.com}) masquerading as a PDF file. VirusTotal identified the payload as malicious, with 46/63 vendors flagging the archive and 60/72 flagging the extracted executable.

VirusTotal Relations linked the sample to multiple phishing emails and ZIP archives submitted throughout 2025, indicating a broad, ongoing distribution campaign.

Hybrid Analysis classified the payload as a Trojan Downloader (Threat Score 100/100). Behavioral observations included:
\begin{itemize}[noitemsep]
    \item Outbound SSL/TLS traffic
    \item TOR-related communication patterns
    \item Multi-process execution
\end{itemize}

No persistence, privilege escalation, or lateral movement behaviors were observed.

\subsection{Assessment}
The phishing email and associated malware reflect low-sophistication cybercrime activity aimed at establishing a foothold on the victim’s system. The combination of spoofed sender identity, public email reply-to address, VPS-based sending infrastructure, and commodity downloader malware aligns with financially motivated threat actors distributing malware at scale.

% ---------------------------------------------------------
% IOC TABLE
% ---------------------------------------------------------
\section{Indicators of Compromise (IOCs)}

\begin{tabularx}{\textwidth}{|l|X|X|}
\hline
\textbf{Category} & \textbf{Indicator} & \textbf{Description} \\
\hline
Sender IP & 192.119.71.157 & Hostwinds VPS; origin of spoofed phishing email \\
\hline
Spoofed Domain & mutawamarine.com & Legitimate business domain impersonated \\
\hline
Reply-To Domain & mail.com & Public email provider used for redirection \\
\hline
Attachment Name & SWT\_ \#09674321\_\_\_\_PDF\_\_.CAB & Malicious RAR/CAB archive disguised as PDF \\
\hline
Extracted Payload & SWT\_ \#09674321\_\_PDF.com & Malicious .NET executable \\
\hline
SHA-256 (Archive) & \seqsplit{2e91c533615a9bb8929ac4bb76707b2444597ce063d84a4b33525e25074fff3f} & Hash of malicious archive \\
\hline
SHA-256 (Payload) & \seqsplit{05261f5a64f81a34fdde66cc82b573773e5dfa3bb5c3ccbfe2d0eef0e9d7b6c9} & Hash of Trojan Downloader \\
\hline
Email Authentication & SPF: Fail; DKIM: None; DMARC: Fail & Confirms spoofing \\
\hline
MX Record & mutawamarine-com.mail.protection.outlook.com & Legitimate Microsoft 365 mail host \\
\hline
Nameservers & DNS1.STABLETRANSIT.COM / DNS2.STABLETRANSIT.COM & Rackspace \\
\hline
Behavioral Indicators & SSL/TLS traffic; TOR activity; Multi-process execution & Downloader behavior \\
\hline
\end{tabularx}

% ---------------------------------------------------------
% RECOMMENDATIONS
% ---------------------------------------------------------
\section{Recommendations}

\subsection{Email Security Controls}
\begin{itemize}
    \item Enforce SPF, DKIM, and DMARC validation.
    \item Quarantine or reject messages failing DMARC alignment.
    \item Block executable content inside compressed archives.
    \item Enable attachment sandboxing.
\end{itemize}

\subsection{User Awareness and Training}
\begin{itemize}
    \item Reinforce phishing awareness around unexpected receipts.
    \item Encourage reporting of suspicious emails.
\end{itemize}

\subsection{Endpoint Protection}
\begin{itemize}
    \item Block unknown or unsigned executables.
    \item Detect downloader behavior and multi-process spawning.
    \item Alert on suspicious outbound encrypted traffic.
\end{itemize}

\subsection{Network Security}
\begin{itemize}
    \item Block TOR exit nodes and known malicious IPs.
    \item Restrict traffic to unused VPS providers.
    \item Enable SSL/TLS inspection where appropriate.
\end{itemize}

\subsection{Incident Response Readiness}
\begin{itemize}
    \item Maintain phishing and malware containment playbooks.
    \item Retain logs for at least 90 days.
\end{itemize}

\subsection{IOC Deployment}
\begin{itemize}
    \item Push IOCs to SIEM, EDR, firewall, and email gateway.
    \item Conduct threat hunting for related activity.
\end{itemize}

% ---------------------------------------------------------
% APPENDIX — METHODOLOGY AND SOURCE
% ---------------------------------------------------------
\section*{Appendix A — Methodology and Source}

This report is based entirely on publicly available information (OSINT).  
No scanning, exploitation, or interaction with live systems occurred.

The phishing email sample (\texttt{challenge.eml}) was obtained from the 
TryHackMe training environment \textit{“The Greenholt Phish”}.  
All analysis was performed on a controlled, sandboxed copy of the file.

\end{document}